\documentclass[11pt,a4paper]{article}

\usepackage[T1]{fontenc}
\usepackage[utf8]{inputenc}
\usepackage{lmodern}
\usepackage{amsmath,amssymb,amsthm}
\usepackage{booktabs}
\usepackage{array}
\usepackage{microtype}
\usepackage{xcolor}
\usepackage[margin=1.05in]{geometry}
\usepackage[numbers,sort&compress]{natbib}
\usepackage[hidelinks,pdfusetitle]{hyperref}
\usepackage{url}
\urlstyle{same}

% ---------------------------------------------------------------------------
% \todo{...}: visible marker for items the author must confirm before submission.
% Set \todofalse to hide them.
\newif\iftodo\todotrue
\newcommand{\todo}[1]{\iftodo{\color{red!80!black}\sffamily\bfseries [TODO: #1]}\fi}
% ---------------------------------------------------------------------------

\theoremstyle{plain}
\newtheorem{theorem}{Theorem}[section]
\newtheorem{lemma}[theorem]{Lemma}
\newtheorem{proposition}[theorem]{Proposition}
\newtheorem{corollary}[theorem]{Corollary}
\theoremstyle{remark}
\newtheorem{remark}[theorem]{Remark}
\newtheorem*{remark*}{Remark}

\newcommand{\Aut}{\operatorname{Aut}}
\newcommand{\Z}{\mathbb{Z}}
\newcommand{\bin}[2]{\binom{#1}{#2}}
\newcommand{\file}[1]{\texttt{#1}}

\title{There is no Leech tree on 18 vertices\\[3pt]
  \large and related distinct-distance tree problems}
\author{Lingsen Meng\\[2pt]
  \small Department of Earth, Planetary, and Space Sciences,\\[-2pt]
  \small University of California, Los Angeles, CA 90095, USA\\[-2pt]
  \small \texttt{meng@epss.ucla.edu}\ \todo{author to confirm e-mail address}}
\date{Draft of 19 August 2026}

\hypersetup{
  pdftitle={There is no Leech tree on 18 vertices, and related distinct-distance tree problems},
  pdfauthor={Lingsen Meng},
  pdfsubject={Leech trees, distinct distance trees, exhaustive search}}

\begin{document}
\maketitle

\begin{abstract}
A \emph{Leech tree} of order $n$ is a tree on $n$ vertices with positive integer
edge weights whose $\bin n2$ pairwise path-weights are exactly $1,2,\dots,\bin n2$.
Leech (1975) found five such trees, of orders $2,3,4,4,6$; Taylor (1977) showed
that the order must be a square or a square plus two; computer searches by
Sz\'ekely, Wang and Zhang (2005) and by Calhoun, Ferland, Lister and Polhill (2007)
excluded the admissible orders $9$, $11$ and $16$, so that $n=18$ has been the
smallest open order.  We show that there is no Leech tree on $18$ vertices.  The
proof is an exhaustive computer search over ``forced forests'' in the sense of
Calhoun et al., organised so that every weighted forest that can occur as the set
of lightest edges of a Leech tree is generated exactly once up to isomorphism.
The search tree has $59{,}779{,}854{,}336$ nodes and contains no tree on $18$
vertices; its per-level node counts are a mathematical invariant (the number of
non-isomorphic forced forests with a given number of edges) which we report and
which was reproduced bit-for-bit by three independent runs on two architectures
and two compilers, and by an independent clean-room implementation written from
the algorithm description alone.  We also record the same statistics for
$n\le 17$, where the search reproduces the known results, and discuss what would
be needed to turn the computation into a formally checkable certificate.
Variants of the same two engines settle several neighbouring questions: for the
minimal distinct distance trees of Calhoun et al.\ we obtain $M(11)=60$, with
exactly two minimal trees, and $77\le M(12)\le 78$ (their bounds were
$59\le M(11)\le 77$, $69\le M(12)\le 94$); there is no modular Leech tree of
order $9$ or of order $11$, while modular Leech trees of order $5$ do exist,
contrary to a statement quoted in the literature; and there are exactly six
leaf-Leech trees with six leaves, five of which contain an irreducible vertex,
which answers Question 3.1 of Ozen, Wang and Yalman in the affirmative.
\end{abstract}

\noindent\textbf{Keywords.} Leech tree, perfect distance tree, distinct distance
tree, minimal distinct distance tree, modular Leech tree, leaf-Leech tree,
Golomb ruler, Sidon set, exhaustive search, computer-assisted proof.

\noindent\textbf{MSC 2020.} 05C05, 05C78, 05C85, 11B83.

\section{Introduction}\label{sec:intro}

Let $T$ be a tree on $n$ vertices with a weight $w(e)\in\Z_{>0}$ on every edge,
and let $d(x,y)$ denote the weight of the (unique) $x$--$y$ path.  Following
Leech~\cite{Leech1975} we call $(T,w)$ a \emph{Leech tree} (also \emph{perfect
distance tree}, \emph{perfect distance-distinct tree}) if the multiset
$\{d(x,y):\{x,y\}\subseteq V(T)\}$ equals $\{1,2,\dots,N\}$ with $N=\bin n2$.
Leech exhibited the five Leech trees of Table~\ref{tab:known} and asked whether
there are any others.  None has been found since.

\begin{table}[ht]
\centering
\caption{The five known Leech trees \cite{Leech1975}.  Weights are listed as
edge lists $u$--$v$:$w$ on vertices $0,1,\dots$; each labelling is unique up to
isomorphism.}
\label{tab:known}
\small
\begin{tabular}{@{}cllll@{}}
\toprule
$n$ & $N$ & tree & weighted edges & remark\\
\midrule
2 & 1  & $K_2$            & 0--1:1 & \\
3 & 3  & $P_3$            & 0--1:1, 1--2:2 & \\
4 & 6  & $K_{1,3}$ (star) & 0--1:1, 0--2:2, 0--3:4 & \\
4 & 6  & $P_4$ (path)     & 0--1:1, 1--2:3, 2--3:2 & weight 3 must be the middle edge\\
6 & 15 & $B_{2,2}$ (double star) & 0--1:1, 0--2:2, 0--3:5, 3--4:4, 3--5:8 & centres $0,3$ joined by weight $5$\\
\bottomrule
\end{tabular}
\end{table}

Necessary conditions on the order were given by Taylor~\cite{Taylor1977}:
counting odd distances shows that a Leech tree of order $n$ exists only if
$n=k^2$ or $n=k^2+2$ for some integer $k$ (Lemma~\ref{lem:taylor} below).  The
admissible orders are therefore $2,3,4,6,9,11,16,18,25,27,\dots$.  Sz\'ekely,
Wang and Zhang~\cite{SzekelyWangZhang2005} excluded $n=9$ (also excluded by Shen
Lin's search reported by Taylor~\cite{Taylor1991}) and $n=11$ by computer search,
and proved asymptotic upper bounds on the length of a path and on the maximum
degree of a Leech tree.  Calhoun, Ferland, Lister and Polhill~\cite{Calhoun2007}
introduced \emph{distinct distance trees} (all path weights distinct, not
necessarily forming an interval), gave a backtracking algorithm over weighted
forests with a \emph{forced next weight} (their Lemma~2.6 and Algorithm~2.7), and
determined all perfect distance trees of order $n<18$; in particular they
excluded $n=16$.  Their summary states the bound as $n\le 24$, but their
abstract, their Table~1 (where the entry for $n=18$ is only the trivial lower
bound), and every later paper give $n<18$; we treat ``24'' as a misprint.
Further structural results---the non-existence of Leech trees of diameter~$3$
for $n\ge 7$, of certain ``brooms'', the ``leaf-Leech'' transformation, and Luo
and Yu's finiteness result for diameter~$4$ and improved degree bound---are in
\cite{OzenWangYalman2016,Varghese2020,LuoYu2024}; a Leech labelling of general
graphs was introduced in \cite{Eldho2024}.  The order $n=18$ is listed as the
smallest open case, for instance in \cite{Varghese2020,OzenWangYalman2016} and
in the FrontierMath open-problem list~\cite{FrontierMath}.

In this note we settle the case $n=18$.

\begin{theorem}\label{thm:main}
There is no Leech tree on $18$ vertices.
\end{theorem}

The proof is a computer search.  Its mathematical content is small---the forcing
lemma of \cite{SzekelyWangZhang2005,Calhoun2007} and an elementary description of
the automorphism group of a forest with distinct edge weights
(Lemma~\ref{lem:aut}), which yields a search tree in which every relevant
weighted forest appears exactly once up to isomorphism---and its cost is modest
(about $6\times10^{10}$ search nodes, a few CPU-hours).  We nevertheless think
that a careful account is warranted, for three reasons.  First, because $n=18$
has been repeatedly cited as open, the result should be stated with a precise
description of what was computed and how it was checked.  Second, the per-level
node counts of the search are a well-defined combinatorial quantity (the number
of non-isomorphic ``forced forests'' with $k$ edges that fit inside a Leech tree
of order $18$), so that the computation is falsifiable in a much stronger sense
than a bare ``no solution found'': any re-implementation must reproduce eighteen
integers, not one bit.  Third, we want to be explicit about the status of the
result as a computer-assisted theorem, including the parts that are and are not
covered by independent verification (Section~\ref{sec:discussion}).

The rest of the note is organised as follows.  Section~\ref{sec:prelim}
collects the lemmas that are used, marking which are classical and which are
(minor) additions.  Section~\ref{sec:algorithm} describes the algorithm and
proves that it is complete and free of isomorphic duplicates.
Section~\ref{sec:verification} describes the verification strategy.
Section~\ref{sec:results} states the results for $n\le 18$, including the
per-level tables.  Section~\ref{sec:variants} reports what the same two engines
give on three neighbouring problems---minimal distinct distance trees, modular
Leech trees, and leaf-Leech trees---where the literature leaves gaps that a
search of this kind can close.  Section~\ref{sec:discussion} discusses the status
of the proof and what remains open.

\section{Preliminaries}\label{sec:prelim}

Throughout, $n$ is the order, $N=\bin n2$, and ``distance'' means weighted path
length.  Weighted forests are considered up to isomorphism (a bijection of
vertices preserving edges and weights).  We say that a weighted forest is
\emph{distance-distinct} if the distances between pairs of vertices in the same
component are pairwise distinct.  All lemmas below are elementary; we include
proofs of the two that carry the main argument (Lemmas~\ref{lem:forcing} and
\ref{lem:aut}) and sketch the rest.

\begin{lemma}[forcing lemma; {\cite[``Find-Next-Weight'']{SzekelyWangZhang2005}},
{\cite[Lemma 2.6]{Calhoun2007}}]\label{lem:forcing}
Let $(T,w)$ be a Leech tree with edge weights $w_1<w_2<\dots<w_{n-1}$, and for
$1\le k\le n-1$ let $F_{k-1}$ be the forest formed by the edges of weights
$w_1,\dots,w_{k-1}$.  Then $w_k$ equals the least positive integer that is not a
distance within a component of $F_{k-1}$.
\end{lemma}

\begin{proof}
Let $m$ be that least missing integer.  If $w_k<m$ then $w_k$ is already a
distance in $F_{k-1}$, contradicting distinctness.  If $w_k>m$ then $m$ is not an
edge weight, and any path realising $m$ has at least two edges, hence every edge
on it has weight $<m<w_k$ and lies in $F_{k-1}$; so $m$ would be a distance in
$F_{k-1}$, contradiction.  Hence $w_k=m$.
\end{proof}

Consequently $w_1=1$, $w_2=2$, and $w_3\in\{3,4\}$ according as the edges of
weight $1$ and $2$ are non-adjacent or adjacent; a Leech labelling of a given
tree is determined by the \emph{order} in which its edges receive their weights.
Note that Lemma~\ref{lem:forcing} fails for graphs with cycles~\cite{Calhoun2007}.

\begin{lemma}[Taylor's parity condition~\cite{Taylor1977}; cf.\ {\cite[Thm.~1]{SzekelyWangZhang2005}}, {\cite[Thm.~1.3]{Calhoun2007}}]\label{lem:taylor}
If a Leech tree of order $n$ exists, then $n=k^2$ or $n=k^2+2$ for an integer
$k$; the vertex classes $A,B$ at even and odd weighted distance from a fixed
vertex satisfy $|A|\,|B|=\lceil N/2\rceil$.  For $n=18$: $N=153$,
$\{|A|,|B|\}=\{11,7\}$.
\end{lemma}

\begin{proof}[Sketch]
$d(x,y)\equiv d(x,z)+d(y,z)\pmod 2$ in a tree; a distance is odd iff its ends lie
in different classes, so $|A|\,|B|=\lceil N/2\rceil$, and
$(|A|-|B|)^2=n^2-4|A|\,|B|\in\{n,n-2\}$.
\end{proof}

\begin{lemma}[cut identity]\label{lem:cut}
For any weighted tree, $\sum_{x<y}d(x,y)=\sum_e c_e\,w(e)$, where
$c_e=s_e(n-s_e)$ is the number of vertex pairs separated by $e$ and $s_e$ is the
number of vertices on one side of $e$.  For a Leech tree the left side is
$N(N+1)/2$.
\end{lemma}

\begin{lemma}[Golomb and containment bounds]\label{lem:golomb}
In a distance-distinct weighted tree:
\begin{enumerate}
\item[(a)] the vertices of any path with $h$ edges, placed on a line at their
cumulative distances, form a Golomb ruler with $h+1$ marks, so $d(x,y)\ge
G(h+1)$ where $G(m)$ is the length of an optimal Golomb ruler with $m$ marks;
\item[(b)] if the tree is Leech, every path weight is at most $N$, so the (hop)
diameter $D$ satisfies $G(D+1)\le N$;
\item[(c)] (containment) if $s_x,s_y$ denote the number of vertices ``behind''
$x$ and $y$ on the $x$--$y$ path (so that $s_xs_y$ pairs have paths containing
the $x$--$y$ path), then in a Leech tree $d(x,y)\le N+1-s_xs_y$; in particular
$w(e)\le N+1-c_e$ for every edge.
\end{enumerate}
\end{lemma}

The values
\[
G(1),\dots,G(17)=0,1,3,6,11,17,25,34,44,55,72,85,106,127,151,177,199
\]
(OEIS A003022 \cite{OEIS}) are used; $G(m)$ for $m\le 13$ was re-verified by
exhaustive search for this project, and $G(14),\dots,G(17)$ are taken from the
literature (Shearer's tables~\cite{Shearer1990}; the proof of the 19-mark ruler
\cite{Dollas1998}; the distributed.net OGR project~\cite{DistributedNetOGR}).
Since $G(16)=177>153=N$, a Leech tree of order $18$ has hop-diameter at most
$14$.  This strengthens the exact finite form of
\cite[Thm.~2]{SzekelyWangZhang2005} (which gives diameter $\le 15$ at $n=18$).

\begin{lemma}[star-Sidon degree bound]\label{lem:sidon}
Let $v$ be a vertex of degree $d$ in a Leech tree of order $n$ with incident
weights $a_1<\dots<a_d$.  Then the $a_i$ and the $\bin d2$ sums $a_i+a_j$
($i<j$) are $d+\bin d2$ pairwise distinct integers in $[1,N]$; call a set of
positive integers with this property (pairwise sums of \emph{distinct} elements
distinct and different from the elements) \emph{star-Sidon}.  Let $S(d)$ be the
minimum of $a_{d-1}+a_d$ over star-Sidon sets of size $d$; then
$S(d)\le a_{d-1}+a_d\le N+1-b_{(1)}b_{(2)}$, where $b_{(1)}\le b_{(2)}$ are the
two smallest branch sizes at $v$.  Exhaustive computation gives
\begin{center}
\begin{tabular}{@{}l*{11}{r}@{}}
\toprule
$d$    & 2 & 3 & 4  & 5  & 6  & 7  & 8  & 9  & 10  & 11  & 12\\
\midrule
$S(d)$ & 3 & 6 & 11 & 19 & 31 & 43 & 63 & 80 & 110 & 138 & 169\\
\bottomrule
\end{tabular}
\end{center}
and $S$ is non-decreasing in $d$.  Since $S(12)=169>153$, the maximum degree of
a Leech tree of order $18$ is at most $11$ (a degree-$11$ vertex is possible only
with $b_{(1)}b_{(2)}\le 16$).
\end{lemma}

\begin{remark*}[correction]
An earlier version of this lemma asserted that the incident weights form a
Golomb ruler (equivalently a Sidon set in the strong sense) and used
$G(d)+G(d-1)+2$ in place of $S(d)$.  That is false: the Leech condition only
makes $\{a_i\}$ a \emph{weak} Sidon set (sums of two distinct elements
distinct); three-term arithmetic progressions among the $a_i$ are not excluded.
The star $K_{1,7}$ with weights $\{1,2,4,8,14,19,24\}$ has all $28$ distances
distinct, yet $a_6+a_7=43<44=G(7)+G(6)+2$; the correct minima are those in the
table (the two sequences differ from $d=6$ on).  The counterexample was found
during the internal audit of the search code (Section~\ref{sec:verification}).
The maximum-degree consequence at $n=18$ ($\Delta\le 11$) is unaffected: it
also follows from an exhaustive search over star-Sidon sets with all elements
and sums $\le 153$, which finds a set of size $11$ (e.g.\
$\{1,2,4,7,12,23,32,41,56,70,83\}$) and none of size $12$.
Lemma~\ref{lem:sidon} is a finite, explicit version of the idea behind
\cite[Thm.~3]{SzekelyWangZhang2005} and, we believe, of \cite{LuoYu2024}; we
could not access the text of \cite{LuoYu2024} and make no claim of priority.
\todo{obtain Luo--Yu 2024 and compare with Lemma~\ref{lem:sidon}}
\end{remark*}

\begin{lemma}[weight bounds; {\cite[Thms.~2.8, 2.9]{Calhoun2007}}]\label{lem:weights}
In a Leech tree of order $n$ the largest edge weight satisfies
$m_1\le\lfloor(4n-1)^2/48\rfloor$ and the second largest
$m_2\le\lfloor(N-1)/2\rfloor$; more generally the weights of any set of edges
lying on a common path sum to at most $N$.  At $n=18$: $m_1\le 105$,
$m_2\le 76$.
\end{lemma}

\paragraph{What is used where.}
It is worth being precise about which of these statements the proof of
Theorem~\ref{thm:main} depends on.  The forest search of
Section~\ref{sec:algorithm} uses only Lemma~\ref{lem:forcing}, the trivial facts
that all distances are distinct and at most $N$ and that a subforest of a Leech
tree has at most $n$ vertices, the ``common path'' case of
Lemma~\ref{lem:weights} (any two edges lie on a common path, so $w(e)+w(f)\le N$;
this rule was found never to prune anything at $n\le 18$), and
Lemma~\ref{lem:aut} below.  Lemmas~\ref{lem:taylor}--\ref{lem:sidon} and the
rest of Lemma~\ref{lem:weights} are \emph{not} needed for Theorem~\ref{thm:main};
they are used in the independent per-topology engine
(Section~\ref{sec:pertopology}) as pruning rules and as topology filters, and are
recorded here because they were verified in the course of the project and
because Lemma~\ref{lem:sidon} corrects a statement that circulated in our own
notes.  Lemmas~\ref{lem:forcing}, \ref{lem:taylor}, \ref{lem:weights} and
Lemma~\ref{lem:golomb}(a),(b) are from the literature (\ref{lem:golomb}(b) as
applied to the diameter is folklore; \cite[Prop.~1.2]{Calhoun2007} uses it for
paths); Lemma~\ref{lem:cut}, Lemma~\ref{lem:golomb}(c), Lemma~\ref{lem:sidon}
(in the corrected form) and Lemma~\ref{lem:aut} do not seem to appear in the
Leech tree literature, though none of them is deep.

\section{The algorithm}\label{sec:algorithm}

\subsection{Forced forests}\label{sec:forced}

Fix $n$ and $N=\bin n2$.  By Lemma~\ref{lem:forcing}, if $(T,w)$ is a Leech tree
and $t\ge 1$, then the subforest $F_t=\{e\in T:w(e)<t\}$ satisfies:
(i) $F_t$ is distance-distinct with all distances in $[1,N]$;
(ii) $F_t$ has at most $n$ vertices (we regard $F_t$ as a forest without
isolated vertices; the vertex set is the set of endpoints of its edges);
(iii) if $t'$ is the least positive integer that is not a distance of $F_t$, then
either $F_t=T$ (in which case $t'=N+1$), or $T$ has an edge $e$ of weight exactly
$t'$ and $F_{t'+1}=F_t+e$.
Call a weighted forest \emph{forced} (for the target $[1,N]$ and order $n$) if it
arises from the empty forest by repeatedly adding an edge whose weight is the
least missing distance, in one of three ways: \emph{join} (the new edge connects
two existing components), \emph{attach} (one endpoint is a new vertex), or
\emph{new} (both endpoints are new; a fresh single-edge component), never
exceeding $n$ vertices and never repeating a distance or exceeding $N$.  This is
exactly the search of \cite[Alg.~2.7]{Calhoun2007}.  By (iii) every Leech tree
of order $n$ is a forced forest with $n-1$ edges and $n$ vertices; conversely a
forced forest with $n$ vertices and $n-1$ edges is a tree in which every value
$1,\dots,N$ is realised (the least missing value has been pushed to $N+1$),
hence a Leech tree.

The natural DFS enumerates forced forests level by level (level $k$ = $k$
edges).  Its cost is dominated by the last few levels, and without symmetry
reduction the same abstract forest is generated many times, since the endpoints
of single-edge components and the components themselves can be listed in many
orders.  The one idea that makes the search cheap is complete isomorph
rejection, for which the following two lemmas suffice.

\subsection{Isomorph rejection}\label{sec:isomorph}

\begin{lemma}[canonical parent]\label{lem:parent}
In a forced forest the edge weights are pairwise distinct and the weights of
successive edges are strictly increasing; hence removing the edge of maximum
weight from a forced forest with $k+1$ edges yields a forced forest with $k$
edges, and this parent is well defined up to isomorphism.
\end{lemma}

\begin{proof}
Each new weight $t$ is the least missing distance and becomes a distance once
the edge is added, so later weights are larger.  Deleting the last edge returns
to the previous state of the DFS.
\end{proof}

\begin{lemma}[automorphisms of a distinct-weight forest]\label{lem:aut}
Let $P$ be a forest without isolated vertices whose edge weights are pairwise
distinct.  Then the group of weight-preserving automorphisms of $P$ is generated
by the transpositions of the two endpoints of the single-edge components of $P$;
in particular $\Aut(P)\cong\Z_2^{\,s}$, where $s$ is the number of single-edge
components, and it acts on components independently.
\end{lemma}

\begin{proof}
Let $\varphi$ be a weight-preserving automorphism.  Since the weights are
distinct, $\varphi$ fixes every edge setwise.  If two distinct edges $e,f$ share
a vertex $v$, then $\varphi(v)\in e\cap f=\{v\}$, so $v$ is fixed.  In a
component with at least two edges every edge has an endpoint of degree $\ge 2$;
that endpoint is fixed, and then the other endpoint of the (setwise fixed) edge
is fixed too.  So every component with $\ge 2$ edges is fixed pointwise.  A
single-edge component admits exactly the identity and the endpoint swap, and
swaps of different components commute and are independent.
\end{proof}

\begin{proposition}[exactly-once generation]\label{prop:once}
Consider the following restricted DFS.  Vertices are labelled in order of
appearance.  At a node $P$ (a labelled forced forest), the \emph{eligible}
vertices are all vertices of $P$ except the larger-labelled endpoint of each
single-edge component.  Children are: all joins between eligible vertices $x<y$
in different components; all attachments of a new vertex to an eligible vertex
$x$; and, if $P$ has at most $n-2$ vertices, the single ``new'' child.  (Each
child is kept only if it is a forced forest, i.e.\ passes the distance and
vertex tests.)  Then every abstract forced forest with $k$ edges is generated
exactly once at level $k$, for every $k$.
\end{proposition}

\begin{proof}
Induction on $k$.  Assume the level-$k$ nodes are pairwise non-isomorphic and
cover every abstract forced forest with $k$ edges.  Let $G$ be an abstract
forced forest with $k+1$ edges, $e$ its maximum-weight edge and $P=G-e$; by
Lemma~\ref{lem:parent}, $P$ is a forced forest with $k$ edges, so it appears
exactly once as a labelled node $P'$.  Two extensions $P'+(x,y,t)$ and
$P'+(x',y',t)$ (with the same forced weight $t$) are isomorphic iff some
$\varphi\in\Aut(P')$ maps $\{x,y\}$ to $\{x',y'\}$: an isomorphism of the
children fixes the unique maximum-weight edge and restricts to an automorphism
of the parents.  By Lemma~\ref{lem:aut}, $\Aut(P')$ is the product of the
endpoint swaps of the single-edge components, so the orbit of an endpoint set is
obtained by independently swapping the endpoint(s) that lie in single-edge
components.  Choosing the smaller label in each single-edge component therefore
selects exactly one representative per orbit for joins (unordered pairs in
distinct components), for attachments (orbit of $x$), and trivially for the
``new'' child.  Children of non-isomorphic parents are non-isomorphic (their
canonical parents would be isomorphic).
\end{proof}

Proposition~\ref{prop:once} is what allows the search tree to be \emph{counted}:
its level-$k$ node count equals the number of non-isomorphic forced forests with
$k$ edges for the given $(n,N)$.  This makes the computation checkable in the
strong sense discussed in the Introduction.  Because $\Aut$ is recomputed from
the current component structure at each node, no automorphism state is cached;
the appearance-order labelling is used only to pick ``the smaller endpoint of a
two-vertex component''.

\subsection{Implementation and sharding}\label{sec:impl}

The engine (\file{forest\_search.cpp}, C++17, no dependencies) keeps, for each
vertex $x$, the bitset $D_0[x]$ of distances from $x$ within its component, and
the bitset $R$ of realised distances.  Adding an edge $(x,y,t)$ generates the new
distances as shifted unions $D_0[x]+t+D_0[y]$; a child is rejected as soon as a
shifted set meets $R$ or exceeds $N$, and a per-vertex prefilter discards $x$
when already the distances from $x$ to the \emph{new neighbour} alone collide
(these are a subset of the new distances of any join or attachment at $x$).  All
state is undone on return (incremental, no copies).  The only other prune is
$t+w_{\max}\le N$ (Lemma~\ref{lem:weights}, common path), which never fires at
$n\le 18$.\footnote{More precisely, the clean-room implementation of
Section~\ref{sec:cleanroom} found that this rule fires exactly once at $n=12$
(removing one level-10 node) and never in the $n=18$ tree; since it is a
necessary condition it cannot remove a solution, and both engines apply it, so
the per-level counts reported here are those \emph{with} the rule.}  There are
no bounds, counting arguments or lookahead in this engine; its efficiency comes
from Proposition~\ref{prop:once} (which alone reduced the node count roughly
tenfold at $n=12$--$13$ relative to the same DFS without isomorph rejection) and
from bit-parallel distance sets.

For parallelism the DFS is sharded at a fixed level $L$: the level-$L$ nodes are
enumerated in DFS order (deterministic), and shard $i$ of $K$ explores the
subtrees of the level-$L$ nodes with index $\equiv i \pmod K$; the levels $\le L$
are visited by every shard.  Hence, if $h_i(d)$ is the number of level-$d$ nodes
visited by shard $i$, one must have $h_i(d)=h(d)$ for $d\le L$ and every $i$,
and $\sum_i h_i(d)=h(d)$ for $d>L$; the total number of visited nodes is
$\sum_i\mathrm{nodes}_i=\mathrm{unique}+(K-1)\cdot\mathrm{prefix}(\le L)$.
These identities were checked in $36$ $(n,\text{target},K,L)$ configurations,
including $K$ not dividing the number of level-$L$ nodes and $K$ larger than it.
The verdict script for a sharded run asserts: exactly $K$ records, indices
$0,\dots,K-1$ each once, identical $(K,L)$, all \texttt{DONE}, and the histogram
identities above.

\section{Verification}\label{sec:verification}

The engine of Section~\ref{sec:algorithm} was written by an AI coding agent
under human direction (see the statement in Section~\ref{sec:ai}).  We
therefore treated it as untrusted and audited it in the following ways.
Everything in this section refers to code that is available with the paper
(Section~\ref{sec:data}).

\subsection{Independent enumerator (oracle) and differential tests}\label{sec:oracle}

A Python enumerator (\file{referee\_forest\_enum.py}), written independently of
the C++ code from the definition in Section~\ref{sec:forced}, generates
\emph{all} labelled join/attach/new children of every node without any symmetry
rule, keeps a child iff its full distance multiset (recomputed from scratch by
BFS) is duplicate-free and inside the target, and de-duplicates each level by a
complete canonical form for distinct-weight forests without isolated vertices
(the sorted multiset of per-vertex sorted incident-weight tuples).  Its
per-level counts equal the engine's for every level and every $n\le 12$ (at
$n=12$: $1,1,2,8,41,229,1384,8877,58933,278539,356479$ nodes on levels
$0,\dots,10$, 23 minutes in Python), for planted targets with holes at
$n=6,\dots,9$, and for the prefix of the $n=17$ and $n=18$ trees through level
$9$ (at $n=18$: $480{,}085$ level-9 forests).  A second Python mode that
re-implements the engine's own generation rule \emph{without} de-duplication
produced zero canonical-form collisions in all these cases, i.e.\ the isomorph
rejection is irredundant as implemented and not only as proved.

Two independent formulations of the same problem were used as further oracles:
(a) a per-topology engine (Section~\ref{sec:pertopology}) whose solution counts,
summed over \emph{all} unlabelled trees of a given order, must equal the forest
engine's; this was checked on $45$ instances ($n=6,\dots,12$, standard and
planted targets, including targets with two non-isomorphic solutions) with $0$
mismatches; and (b) a plain edge-by-edge backtracking over target values with
\emph{no} forcing lemma, summed over topologies as labellings divided by
$|\Aut|$, for $n\le 9$.  Full-depth planted instances at $n=13,\dots,18$
(random weighted trees with distinct distances, custom target set) were run
through the engine: in all $16$ instances the planted tree was found among the
solutions and every reported solution was re-checked; these are the only tests
that exercise the $n=18$, full-vertex-count code path with a witness.  Finally
the production binary was shown to be a faithful build of the archived source by
reproducing its depth histograms at $n=4,\dots,14$ and the exact node count
($100{,}024{,}625$) of one $n=18$ production shard from a byte-for-byte rebuild.

\subsection{Independent method: per-topology search}\label{sec:pertopology}

Before the forest engine we built a per-topology exact search
(\file{leech\_search.cpp}): for each of the $123{,}867$ unlabelled trees on $18$
vertices (OEIS A000055 \cite{OEIS}; enumerated by two independent generators,
the WROM algorithm \cite{WROM1986} as implemented in NetworkX
\cite{HagbergNetworkX2008} and nauty's \texttt{gentreeg} \cite{McKayPiperno2014},
which agree for $n=5,9,11,16,18$: $3,47,235,19320,123867$) it branches on which
unassigned edge receives the least missing value (Lemma~\ref{lem:forcing}) and
prunes with distinctness, the containment and Golomb windows
(Lemma~\ref{lem:golomb}), the cut identity (Lemma~\ref{lem:cut}), Hall-type
counting on value windows, Taylor's parity (Lemma~\ref{lem:taylor}),
Lemma~\ref{lem:weights}, and complete symmetry breaking over $\Aut$ of the
topology (verified as \emph{exactly one representative per orbit} by the
identity $\mathrm{count(sym)}\cdot|\Aut|=\mathrm{count(nosym)}$ on planted
instances with $|\Aut|$ up to $20{,}000$).  This engine was audited by a
differential harness against a rule-free DFS and against a CP-SAT enumeration
\cite{ORTools}: $574$ instances (all topologies $n=7,\dots,10$ with planted
targets, the standard target $n=7,\dots,9$, parity-skewed families, symmetry
stress set) $\times$ $3$ binaries $\times$ $32$ flag subsets, and a second seed
with $560$ instances, with $0$ discrepancies.  It reproduces the literature
id-by-id at $n=9$ ($47$ topologies) and $n=11$ ($235$), agreeing with CP-SAT and
with the DRAT-certified SAT route (Section~\ref{sec:sat}), and it has now decided
\emph{all} $19{,}320$ topologies of order $16$ on the cluster as UNSAT
($19{,}308$ with the first-generation engine at a $1800$\,s cap; the $12$
topologies that hit the cap were re-run with the improved engine and finished in
$673$--$1317$\,s each).  This is an independent-method reproduction of the
non-existence of a Leech tree of order $16$ \cite{Calhoun2007}.

The per-topology engine is between three and four orders of magnitude slower in
total than the forest engine (it re-derives the same small-weight forests once
per topology: $n=16$ took $321$ CPU-s with the forest engine against roughly
$500$--$700$ CPU-h topology-by-topology), which is why it is used as a
\emph{cross-check} rather than as the primary proof.

At $n=18$ this cross-check is \emph{incomplete}, and we state its status
precisely.  A first pass over all $122{,}344$ surviving topologies (those left
after the proven filters of Section~\ref{sec:order18}) has finished on the
cluster with a $300$\,s time cap per topology: $73{,}021$ topologies UNSAT,
$49{,}323$ undecided at the cap, and---this is the part that matters---$0$
topologies SAT.  The $49{,}323$ undecided topologies are being re-run with the
improved engine at a $3600$\,s cap; that run had not finished at the time of
writing.  The independent-method cross-check therefore currently covers
$73{,}021$ of the $122{,}344$ survivor topologies, i.e.\ $59.7\%$ of them; the
remaining $1{,}523$ of the $123{,}867$ topologies had been removed beforehand by
the filters of Section~\ref{sec:order18}, and the forest engine of
Theorem~\ref{thm:main} uses no topology filters at all.  No topology anywhere in
the project, at any order, has ever been reported SAT except for those carrying
the three known Leech trees of orders $4$ and $6$.  Until the
re-run is complete, Theorem~\ref{thm:main} rests on the forest engine, its
audits (Section~\ref{sec:oracle}), the three bit-identical replications
(Section~\ref{sec:replication}) and the clean-room implementation
(Section~\ref{sec:cleanroom}); the per-topology run is corroboration by a
different method, not a second complete proof.
\todo{report the completed per-topology outcome for $n=18$ once the re-run of
the $49{,}323$ undecided topologies has finished; if it disagrees with the
forest engine, the theorem must be withheld.}

\subsection{SAT encoding with checked proofs ($n\le 11$)}\label{sec:sat}

For small orders we also produced machine-checkable certificates: an
order-encoding CNF per topology (pair/value variables, ternary sum relations
along a rooted tree, sibling-subtree symmetry breaking), solved with CaDiCaL
\cite{BiereCaDiCaL2020,BiereCaDiCaL2024} (version 3.0.1; Kissat 4.0.4 was used
as a second solver) producing DRAT proofs, checked with \texttt{drat-trim}
\cite{WetzlerHeuleHunt2014}.  All $47$ topologies at $n=9$ and all $235$ at
$n=11$ are UNSAT with verified proofs (also the $5$ non-Leech topologies at
$n=6$; the sixth yields the known tree); this reproduces
\cite{SzekelyWangZhang2005} with independent certificates.  The approach does
not scale ($30$ random $n=16$ topologies: $0/30$ decided in $600$\,s each), so
no DRAT-style certificate exists for $n=18$; see Section~\ref{sec:discussion}.

\subsection{Replication of the $n=18$ run}\label{sec:replication}

The $n=18$ forest search was run three times with the reference engine, with
two shard layouts, on two architectures and with two compilers, and once more
with the clean-room implementation of Section~\ref{sec:cleanroom}; the
per-shard depth histograms were compared (Table~\ref{tab:replication}).

\begin{table}[ht]
\centering
\caption{Replication of the $n=18$ forest search.  $\sum$nodes counts the shared
prefix (levels $\le L$) once per shard; ``unique'' $=\sum\text{nodes}-(K-1)\cdot
73{,}408$, where $73{,}408$ is the number of nodes on levels $0,\dots,8$.
Hoffman2 is the UCLA shared cluster (SLURM array jobs); ``local'' is an Apple
M4 laptop.}
\label{tab:replication}
\footnotesize
\setlength{\tabcolsep}{3.5pt}
\begin{tabular}{@{}lllrrlcc@{}}
\toprule
run & machine / compiler & $(K,L)$ & $\sum$nodes & unique nodes & levels 9--16 & L17 & trees\\
\midrule
Hoffman2 job 83703 & x86\_64, gcc 11.5 & (210, 8) & 59,795,196,608 & 59,779,854,336 & Table~\ref{tab:levels18} & 0 & 0\\
Hoffman2 job 83752 & x86\_64, gcc 11.5 & (175, 9) & 59,876,162,118 & 59,779,854,336 & identical & 0 & 0\\
local, reference & arm64 M4, clang & (512, 8) & 59,817,365,824 & 59,779,854,336 & n/a\footnotemark & --- & 0\\
local, clean-room & arm64 M4, clang & (100, 8) & 59,787,121,728 & 59,779,854,336 & identical & 0 & 0\\
\bottomrule
\end{tabular}
\end{table}
\footnotetext{The local driver of the reference engine discarded the per-shard
stderr histograms; only the node counts and status of each shard were kept.
The identity of the unique node count with the two cluster runs is the check
available for this run.}

The four sums of visited nodes differ exactly by the multiplicity of the shared
prefix ($(K-1)\cdot\mathrm{prefix}(\le L)$), and the unique node counts agree
bit-for-bit.  For the two cluster runs and the clean-room run the per-level
sums agree with each other on every level, and agree with the Python oracle on
levels $\le 9$.  The node counts of the reference engine are deterministic, so
its three-fold replication detects hardware, compiler and job-management faults
(lost or truncated shards, mis-sharding), not algorithmic errors; the
algorithmic checks are those of Sections~\ref{sec:oracle}--\ref{sec:pertopology}
and \ref{sec:cleanroom}.

\subsection{Clean-room re-implementation}\label{sec:cleanroom}

A second implementation of the forced-forest search (\file{cr\_forest.c}, C,
single file) was written from the algorithmic description of
Section~\ref{sec:algorithm} and \cite{Calhoun2007} only, by an agent that had
no access to the production source; it uses its own state representation
(state copied per DFS level, full within-component distance matrix, component
identifiers by smallest label) and its own ordering of children, and shares
with the reference engine only the mathematical rules of
Proposition~\ref{prop:once}.  It comes with its own brute-force Python oracle
(all labelled children, BFS re-computation of distances, canonical-form
de-duplication per level), with which it agrees at every level for $n=4,\dots,12$,
and it passes the same sharding identities on $11$ $(n,K,L)$ cases.  It
reproduces the reference engine's per-level counts at every level for
$n=4,\dots,16$ (at $n=16$: $1{,}096{,}039{,}152$ nodes, identical histogram),
for levels $0,\dots,12$ at $n=17$, and---run over the full $n=18$ tree with
$K=100$, $L=8$ ($34{,}207$ CPU-s at low priority on the laptop)---all eighteen
per-level counts of Table~\ref{tab:levels18}, the unique node count
$59{,}779{,}854{,}336$, and no tree on $18$ vertices.  The verdict script
\file{cleanroom/verify\_18.py} performs the shard-integrity checks of
Section~\ref{sec:impl} and prints the comparison with the reference counts.

\section{Results}\label{sec:results}

\subsection{Orders $n\le 17$}\label{sec:small}

The forest engine finds exactly the known Leech trees and no others: at $n=4$
two (the star and the path), at $n=6$ one (the double star), and none at $n=5$
and $7\le n\le 17$.  (The condition of Lemma~\ref{lem:taylor} is not imposed, so
the non-admissible orders are searched as well and come out empty, as they
must.)  Node counts and CPU times (single core, Apple M4, under load) are given
in Table~\ref{tab:small}.

\begin{table}[ht]
\centering
\caption{Forest search for $n\le 17$ (target $[1,\bin n2]$).  ``Last levels''
lists the node counts of the three deepest non-empty levels.  For $n=17$ the
run was sharded ($K=32$, $L=8$); the unique node count subtracts the $31$
repeated copies of the $73{,}408$-node prefix.}
\label{tab:small}
\small
\begin{tabular}{@{}rrrcl@{}}
\toprule
$n$ & nodes (unique) & CPU & Leech trees & last levels\\
\midrule
11 & 123,309 & 0.05\,s & 0 & L7 8,645; L8 43,855; L9 69,144\\
12 & 704,494 & 0.14\,s & 0 & L8 58,933; L9 278,539; L10 356,479\\
13 & 4,178,290 & 0.9\,s & 0 & L9 424,878; L10 1.86M; L11 1.82M\\
14 & 25,486,327 & 6.4\,s & 0 & L10 3.25M; L11 13.0M; L12 8.66M\\
15 & 162,497,458 & 44\,s & 0 & L11 26.2M; L12 93.5M; L13 38.4M\\
16 & 1,096,039,152 & 321\,s & 0 & L12 220.0M; L13 683.5M; L14 154.7M\\
17 & 7,875,306,893 & 1.7 CPU-h & 0 & (histograms not retained)\\
\bottomrule
\end{tabular}
\end{table}

The full level histogram at $n=16$ (levels $0,\dots,15$) is
\begin{multline*}
1,\ 1,\ 2,\ 8,\ 41,\ 229,\ 1384,\ 8899,\ 62843,\ 480048,\ 3968025,\ 33269745,\\
220001476,\ 683511193,\ 154735257,\ 0.
\end{multline*}  The counts at $n=16$ and $n=17$ reproduce
the non-existence results of \cite{Calhoun2007}; those at $n=9,11$ reproduce
\cite{SzekelyWangZhang2005}, and are also confirmed by the per-topology engine,
CP-SAT, and (for $n=9,11$) DRAT-checked SAT proofs.  The per-level counts at
low levels are the same for all $n$ once $n$ is large enough for the vertex
bound not to bite ($1,1,2,8,41,229,1384,8899,62843,\dots$), and grow by a factor
of roughly $7.5$ per level; the whole tree grows by a factor of $6.5$--$7$ per
unit of $n$.

\subsection{Order 18}\label{sec:order18}

\begin{theorem}[Theorem~\ref{thm:main}, quantitative form]\label{thm:quant}
For $n=18$, $N=153$, the search tree of Section~\ref{sec:algorithm} has the
per-level node counts of Table~\ref{tab:levels18}, a total of
$59{,}779{,}854{,}336$ nodes, and no node at level $17$.  Consequently there is
no Leech tree of order $18$.
\end{theorem}

\begin{table}[ht]
\centering
\caption{Number of non-isomorphic forced forests with $k$ edges for
$(n,N)=(18,153)$.  Verified identical in both cluster runs and in the
clean-room implementation; levels $\le 9$ also by the Python oracle.}
\label{tab:levels18}
\small
\begin{tabular}{@{}l*{6}{r}@{}}
\toprule
level $k$ & 0 & 1 & 2 & 3 & 4 & 5\\
forests & 1 & 1 & 2 & 8 & 41 & 229\\
\midrule
level $k$ & 6 & 7 & 8 & 9 & 10 & 11\\
forests & 1,384 & 8,899 & 62,843 & 480,085 & 3,984,162 & 35,540,837\\
\midrule
level $k$ & 12 & 13 & 14 & 15 & 16 & 17\\
forests & 332,597,341 & 2,896,330,052 & 17,017,193,146 & 37,669,242,243 & 1,824,413,062 & \textbf{0}\\
\bottomrule
\end{tabular}
\end{table}

The tree dies at the last two levels: of the $1.8\times10^9$ distance-distinct
forced forests with $16$ edges (i.e.\ one edge short of a spanning tree on $18$
vertices, or $17$ edges on fewer vertices), none can be completed.  The whole
computation cost about $5$--$11$ CPU-hours depending on machine and load
(roughly $0.3\,\mu$s per node unloaded; $39{,}708$ CPU-s for the local
$512$-shard run at low priority); on the cluster each of the $210$ shard
processes ran for about $1.5$ minutes of wall time.  For comparison, a
per-topology search of the same order was projected at $10^4$--$10^5$ CPU-hours.

We record for completeness that the published structural results prune very
little at the topology level: of the $123{,}867$ trees on $18$ vertices, the
diameter bound (Lemma~\ref{lem:golomb}(b)), the degree bound
(Lemma~\ref{lem:sidon}), the exclusion of stars, double stars, paths and the
broom of \cite{Varghese2020} remove $152$, leaving $123{,}715$; adding real
relaxations (majorisation of the distance vector, projection bound) leaves
$122{,}344$.  The forest search does not use topologies at all.

\section{Neighbouring problems settled by the same engines}\label{sec:variants}

The forced-forest search of Section~\ref{sec:algorithm} and the per-topology
search of Section~\ref{sec:pertopology} are not specific to the target
$\{1,\dots,N\}$.  With small modifications they decide three neighbouring
problems on which the literature stops short of exact values, and on which
the published bounds are, in each case, more than twenty years old or explicitly
left open.  All results in this section were obtained after the $n=18$
computation, with copies of the engines kept separate from the production
sources; every witness reported below was re-checked by an independent
breadth-first-search checker that shares no code with any engine, and the ones
printed here are small enough to be verified by hand.

\subsection{Minimal distinct distance trees: $M(11)=60$ and $77\le M(12)\le 78$}
\label{sec:mdd}

Calhoun, Ferland, Lister and Polhill~\cite{Calhoun2007} define: ``If all of
these distances are distinct, we call the tree a distinct distance tree.  Define
the function $M(n)$ to be the smallest integer so that there exists a distinct
distance tree with $n$ vertices and maximum distance $M(n)$'' \cite[p.~33]{Calhoun2007};
a \emph{minimal distinct distance tree} is one attaining $M(n)$, and a Leech tree
is exactly a distinct distance tree with $M(T)=\bin n2$.  They determine
$M(n)$ for $n\le 10$ ($0,1,3,6,11,15,22,30,39,50$) and give the bounds of their
Table~1 for $11\le n\le 18$, among them $59\le M(11)\le 77$, $69\le M(12)\le 94$
and $80\le M(13)\le 119$, with the remark that ``it will take too long for the
algorithm to check higher numbers of gaps and find the exact value of $M(n)$''
\cite[p.~44]{Calhoun2007}.

The forest engine adapts to $M(n)$ by replacing the forcing lemma with a
three-way decision at each value level $t$ (all values $<t$ decided): either $t$
is already a distance, or $t$ is declared a \emph{gap} (never a distance; at most
$D-\bin n2$ gaps are available when the target maximum is $D$), or $t$ is the
weight of the next edge, added by join, attach or new as before.  Every state
$(\text{forest},t)$ is again reached by exactly one decision sequence, so the
isomorph rejection of Proposition~\ref{prop:once} carries over verbatim and the
number of solutions at $D=M(n)$ is the number of minimal trees up to isomorphism.
Feasibility is decided for increasing $D$; the run at $D=M(n)-1$ is the
exclusion.  The engine reproduces all of Calhoun's values and, where they are
stated, their counts of minimal trees ($n=7$: $2$; $n=8$: $2$; $n=9$: $1$;
$n=10$: $1$, with the same weights $1,2,3,4,5,11,14,18,28$), at costs from $215$
nodes ($n=7$, $D=21$) to $2.2\times10^7$ nodes ($n=10$, $D=49$).

\begin{theorem}\label{thm:mdd}
$M(11)=60$, and there are exactly two minimal distinct distance trees on $11$
vertices, namely (edge lists $u$--$v$:$w$)
\begin{align*}
&0\text{--}1{:}1,\ 0\text{--}2{:}2,\ 3\text{--}4{:}4,\ 5\text{--}6{:}5,\ 5\text{--}7{:}6,\
3\text{--}8{:}8,\ 3\text{--}5{:}9,\ 3\text{--}9{:}16,\ 2\text{--}7{:}26,\ 7\text{--}10{:}27,\\
&0\text{--}1{:}1,\ 0\text{--}2{:}2,\ 3\text{--}4{:}4,\ 5\text{--}6{:}5,\ 5\text{--}7{:}6,\
3\text{--}8{:}8,\ 3\text{--}5{:}9,\ 3\text{--}9{:}16,\ 7\text{--}10{:}26,\ 0\text{--}7{:}27,
\end{align*}
whose $55$ distances are $\{1,\dots,60\}$ minus $\{7,10,21,36,54\}$ and
$\{7,10,21,36,56\}$ respectively.  Moreover $77\le M(12)\le 78$ and
$M(13)\le 105$.
\end{theorem}

The computations behind Theorem~\ref{thm:mdd} are as follows (all on the cluster
unless stated; ``UNSAT'' means the exhaustive run at that $D$ found no tree and
all shards completed).  At $n=11$: $D=59$ UNSAT ($1.80\times10^8$ nodes locally,
$1.83\times10^8$ in $7$ shards on the cluster), $D=60$ two trees
($5.86\times10^8$ nodes, $111$ CPU-s).  For $D=61,\dots,65$ the numbers of
distinct distance trees with maximum distance exactly $D$ are $0,6,3,12,32$.
As an independent-method check, a per-topology backtracking program with no
forcing lemma, no forest search and no isomorph rejection was run over all
$235$ topologies of order $11$: at $D=59$ it finds $0$ labellings in
$12{,}620{,}879{,}225$ nodes and at $D=60$ it finds $18$ labellings in
$15{,}106{,}490{,}676$ nodes, with identical node counts and identical solution
lists on two machines (arm64/clang and x86\_64/gcc); the $18$ labellings are
exactly the two trees above counted with their automorphisms ($6$ and $12$).
The same check at $n=10$ gives $0$ labellings at $D=49$ and $6$ (one tree,
$|\Aut|=6$) at $D=50$.  At $n=12$ the exhaustive runs at
$D=69,70,\dots,76$ are all UNSAT, at costs $3.7\times10^8$, $1.5\times10^9$,
$5.2\times10^9$, $1.6\times10^{10}$, $4.2\times10^{10}$, $1.0\times10^{11}$,
$2.4\times10^{11}$ and $5.1\times10^{11}$ nodes, so $M(12)\ge 77$; and the
witness-hunting mode found the tree
\[
0\text{--}1{:}1,\ 0\text{--}2{:}2,\ 3\text{--}4{:}4,\ 5\text{--}6{:}5,\ 0\text{--}3{:}6,\
5\text{--}7{:}13,\ 8\text{--}9{:}14,\ 4\text{--}8{:}15,\ 6\text{--}9{:}16,\ 8\text{--}10{:}17,\
9\text{--}11{:}37
\]
with all $66$ distances distinct and maximum $78$, so $M(12)\le 78$.  The
exhaustive run at $D=77$ (expected $\sim10^{12}$ nodes) is queued; a search for a
witness at $D=77$ was still running, without success, after $5\times10^{10}$
nodes.  \todo{state the exact value of $M(12)$ once the exhaustive $D=77$ run
completes}  Appending one leaf to the $n=12$ witness with $D=79$ gives a
$13$-vertex distinct distance tree of maximum distance $105$, whence
$M(13)\le 105$ (Calhoun: $119$).

\subsection{Modular Leech trees: none of order 9 or 11, but two of order 5}
\label{sec:modular}

Leach and Walsh introduced the modular relaxation, and Leach~\cite{Leach2014}
states it as: ``Let $T$ be a tree on $n$ vertices and let $k=\bin n2+1$.  We say
that $T$ is a modular Leech tree if there exists an edge weighting function
$w:E(T)\to\Z_k$ such that each of the $\bin n2$ paths within $T$ has a distinct
weight from $1,2,\dots,\bin n2$ with the sums taken modulo $\bin n2+1$''
\cite[p.~1]{Leach2014}; equivalently $w$ induces a bijection between the paths of
$T$ and $\Z_k\setminus\{0\}$.  Every Leech tree is a modular Leech tree, so
orders $2,3,4,6$ occur.  Leach~\cite[Thms.~1--2]{Leach2014} shows that if
$n\equiv 2,3\pmod 4$ then a modular Leech tree of order $n$ forces $n=m^2+2$
(so none exist for $n=7,10,14,15,\dots$), that multiplying a labelling by a unit
of $\Z_k$ gives a labelling, and, by computer search, that there is exactly one
modular Leech tree of order $8$ (over $\Z_{29}$, unique up to group and graph
isomorphism).  Nothing is stated there about $n\ge 9$.

Our engine here is a per-topology depth-first search over the frozen topology
lists, placing edges in breadth-first order from a maximum-degree root and
maintaining the set of realised residues in a $128$-bit word; solutions are
counted as orbits under the unit action of $\Z_k$.  Two implementations with
different symmetry breaking (restriction of the first label versus a
lexicographic-minimum test over the unit stabiliser, the latter with forward
checking) were written and give identical orbit counts for $n=4,\dots,9$;
brute force over all $(k-1)^{n-1}$ labellings independently confirms $n=4,5,6$;
and at $n=8$ the search reproduces Leach exactly, finding the single topology of
his Figure~3 with one labelling up to the leaf swap, our labelling being
$3$ times his in $\Z_{29}$.

\begin{table}[ht]
\centering
\caption{Modular Leech trees of order $n$ over $\Z_k$, $k=\bin n2+1$.  ``Orbits''
counts $\Z_k$-Leech labellings up to multiplication by a unit (graph
automorphisms are not divided out).  New results in bold.}
\label{tab:modular}
\small
\begin{tabular}{@{}rrrlll@{}}
\toprule
$n$ & $k$ & topologies & trees / orbits & search nodes & remark\\
\midrule
4  & 7  & 2   & 2 / 4          & $16$--$18$ & Leech's two trees\\
5  & 11 & 3   & \textbf{2 / 4} & $58$--$127$ & see the discussion below\\
6  & 16 & 6   & 2 / 34         & $1.7\times10^3$--$2.2\times10^4$ & the Leech tree and one other\\
7  & 22 & 11  & 0              & $3.1\times10^4$--$4.6\times10^5$ & agrees with \cite[Thm.~2]{Leach2014}\\
8  & 29 & 23  & 1 / 2          & $2.4\times10^5$--$7.6\times10^5$ & $=$ \cite{Leach2014}\\
9  & 37 & 47  & \textbf{0}     & $6.5\times10^6$--$1.8\times10^7$ & new\\
10 & 46 & 106 & 0              & $3.6\times10^8$--$1.4\times10^{10}$ & consistent with \cite[Thm.~1]{Leach2014}\\
11 & 56 & 235 & \textbf{0}     & $1.3\times10^{10}$ & new; single implementation\\
\bottomrule
\end{tabular}
\end{table}

Table~\ref{tab:modular} summarises the outcome.  The two new non-existence
results are $n=9$ and $n=11$; both orders are permitted by Leach's theorems
($9=3^2$, $11=3^2+2$), so they were open.  The order-$9$ result is produced by
both implementations; the order-$11$ result comes from the faster implementation
only (the slower one did not finish), and we label it accordingly as a
single-implementation computation.

The order-$5$ entry requires comment.  Leach writes: ``By Leach and Walsh [6] and
Theorem 2 we know that none exist for 5 or 7'' \cite[p.~2]{Leach2014}, his~[6]
being \cite{LeachWalsh2011}.  His
Theorems~1--2 exclude $n=7$; they say nothing about $n=5$ (as $5\equiv1\pmod4$),
so the order-$5$ statement rests entirely on \cite{LeachWalsh2011}, which we were
not able to obtain.  Under the definition quoted above, modular Leech trees of
order $5$ do exist.  There are two, on two different topologies; over
$\Z_{11}$ they are the path
$0\text{--}1{:}5,\ 1\text{--}2{:}1,\ 2\text{--}3{:}2,\ 3\text{--}4{:}7$
and the spider
$0\text{--}1{:}1,\ 1\text{--}2{:}5,\ 0\text{--}3{:}3,\ 0\text{--}4{:}7$,
with $4$ unit-orbits in total.  The ten path sums of the first are
$5,1,2,7,6,3,9,8,10$ and $15\equiv 4$, and those of the second are
$1,5,6,3,7,4,8,9,13\equiv 2$ and $10$: in both cases exactly
$\Z_{11}\setminus\{0\}$, as the reader can check in a minute.  Both labellings
use distinct nonzero labels, so they also satisfy the more restrictive readings
of the definition.  We found these by exhaustive search over all
$10^4$ labellings of all three topologies as well as by both search engines.  We
therefore report a discrepancy with the claim quoted above, without being able to
locate its source: it may be a difference of definition in \cite{LeachWalsh2011},
which we have not read, or an error in that paper or in its summary.
Everything else we compute agrees with \cite{Leach2014}.

\subsection{Leaf-Leech trees: six leaves, and Question 3.1 of Ozen--Wang--Yalman}
\label{sec:leafleech}

Ozen, Wang and Yalman~\cite{OzenWangYalman2016} define: ``A (not weighted) tree
$T$ with $n$ leaves is a leaf-Leech tree if the distances between pairs of leaves
are exactly $3,4,\dots,\bin n2+2$'' \cite[Def.~1]{OzenWangYalman2016} (distances
$1$ and $2$ between leaves are impossible except in trivial cases, whence the
shift).  They prove the Taylor analogue (the number of leaves is $k^2$ or
$k^2+2$), that the \emph{expansion} $T^e$ of a Leech tree---subdivide each edge
into a path of its weight and hang a pendant vertex at every original
vertex---is leaf-Leech, and, conversely, that a leaf-Leech tree with no
\emph{irreducible} vertex (a vertex ``with degree at least three and no leaf
neighbors'' \cite[p.~4]{OzenWangYalman2016}) is the expansion of a Leech tree.
They then ask \cite[Question 3.1]{OzenWangYalman2016}: ``Does there exist a
leaf-Leech tree that contains an irreducible vertex?'', remarking that they have
not been able to find a leaf-Leech tree that is not the expansion of a Leech
tree; before the present note the largest number of leaves for which any
leaf-Leech tree was known was $6$, from Leech's order-$6$ tree.

Contracting the degree-$2$ vertices turns a leaf-Leech tree with $L$ leaves into
a series-reduced tree (all internal degrees $\ge 3$) with $L$ leaves and positive
integer edge weights whose $S=\bin L2$ weighted leaf-to-leaf distances are
exactly $\{3,\dots,S+2\}$; conversely every such weighted tree expands to a
leaf-Leech tree.  The series-reduced topologies ($L+1\le |V|\le 2L-2$) were
enumerated with nauty's \texttt{gentreeg} and cross-checked against NetworkX, and
each was solved with CP-SAT \cite{ORTools} in enumerate-all-solutions mode
(weights in $[1,S+2]$, an \texttt{AllDifferent} over the $S$ leaf-pair distance
variables constrained to $[3,S+2]$, which forces the bijection); solutions were
de-duplicated modulo isomorphism of the expanded unweighted tree and re-checked
by the independent checker.

\begin{proposition}\label{prop:leafleech}
There are exactly six leaf-Leech trees with six leaves.  Exactly one of them, the
expansion of Leech's order-$6$ tree, has no irreducible vertex; the other five
have one.  In particular the answer to Question 3.1 of
\cite{OzenWangYalman2016} is yes.
\end{proposition}

For $L=3$ and $L=4$ the enumeration gives $1$ and $2$ leaf-Leech trees, namely
the expansions of the Leech trees of those orders, and for $L=5$ none (as
Proposition~1 of \cite{OzenWangYalman2016} already excludes).  That exactly one
of the six trees at $L=6$ is irreducible-vertex-free is a consistency check on
the enumeration, since by their Theorem~2 such a tree must be the expansion of a
Leech tree of order $6$, and there is exactly one of those.  An explicit witness
for Proposition~\ref{prop:leafleech}, in contracted form, is
\[
1\text{--}0{:}1,\ 1\text{--}2{:}6,\ 1\text{--}5{:}4,\ 0\text{--}6{:}1,\
0\text{--}9{:}1,\ 2\text{--}3{:}6,\ 2\text{--}4{:}2,\ 6\text{--}7{:}3,\
6\text{--}8{:}1 ,
\]
whose expansion has $26$ vertices and six leaves $3,4,5,7,8,9$ at the fifteen
pairwise distances $3,4,\dots,17$; the vertices $1$ and $2$ have degree $3$ and,
in the expanded tree, no leaf neighbour, so both are irreducible.  This tree is
therefore not the expansion of any Leech tree.

The next admissible numbers of leaves are $L=9$ and $L=11$ ($73$ and $488$
series-reduced topologies), where CP-SAT is no longer adequate---exactly as in
the Leech problem itself, it decides the small cases and then stalls (the
$10$-vertex topology at $L=9$ is infeasible in $41$\,s, the $11$-vertex ones were
undecided after $3600$\,s).  A forcing-style engine over leaf distances,
analogous to Section~\ref{sec:algorithm}, would be the natural next step; we
leave it open.

\subsection{Status of the results of this section}\label{sec:variantstatus}

These are computational results of the same kind as Theorem~\ref{thm:main}, with
weaker but explicitly stated verification.  The existence statements
(the two minimal trees at $n=11$, the $D=78$ tree at $n=12$, the modular trees of
order $5$, the six leaf-Leech trees) are witnesses, printed above or in the
archive, and are checkable by hand or by any independent checker; they do not
depend on the search being exhaustive.  The non-existence statements
($M(11)\ge 60$, $M(12)\ge 77$, no modular Leech tree of order $9$ or $11$, only
six leaf-Leech trees at $L=6$) do depend on exhaustiveness: $M(11)\ge 60$ and the
count of minimal trees at $n=11$ are confirmed by a second, structurally
different program on two architectures; the modular order-$9$ result by two
implementations with different symmetry breaking; the modular order-$11$ result
by one implementation only; the leaf-Leech enumeration by a single CP-SAT model
over an independently cross-checked topology list.  We label them accordingly and
regard the modular order-$11$ result as the least independently supported of the
group.

\section{Discussion}\label{sec:discussion}

\paragraph{Status of the result.}
Theorem~\ref{thm:main} is a computer-assisted theorem in the usual sense: the
mathematics (Lemmas~\ref{lem:forcing}, \ref{lem:parent}, \ref{lem:aut} and
Proposition~\ref{prop:once}) is short and has been checked by hand and by
exhaustive small cases; the remainder is a finite computation whose correctness
depends on the faithful execution of a program.  What distinguishes this
computation from a bare ``no solution found'' is that its intermediate
output---the eighteen per-level counts of Table~\ref{tab:levels18}---is a
mathematically defined sequence (the number of non-isomorphic forced forests of
each size), so that any independent implementation can be checked against it
level by level.  The shallow levels have been confirmed by an oracle written
from the definitions; all eighteen levels have been reproduced by an independent
implementation written from the algorithm description; and three replications
of the reference engine on different hardware and compilers agree bit-for-bit.
The one verification we would like to have and do not yet have in full is the
independent \emph{method}: the per-topology engine, which shares no code and no
search order with the forest engine, has decided $73{,}021$ of the $122{,}344$
surviving $18$-vertex topologies (all UNSAT, none SAT) and is still running on
the remainder (Section~\ref{sec:pertopology}).  It has, however, already
reproduced the complete order-$16$ result ($19{,}320$ topologies, all UNSAT) by
this independent route, which is the strongest available calibration of the two
engines against each other and against \cite{Calhoun2007}.

\paragraph{Certificates.}
We do not have a proof certificate for $n=18$ that could be checked by a
formally verified checker, as we do for $n\le 11$ (DRAT).  Two routes are open.
(i) Cube-and-conquer: use the forcing-lemma tree to a fixed depth as cubes and
certify each cube with a CDCL solver producing DRAT proofs; the completeness of
the cube set follows from Lemma~\ref{lem:forcing} by a one-paragraph argument,
so the search engine's pruning need not be trusted, only its branching order.
Our experiments at $n=16$ suggest a cost of the order of one CPU-hour per
topology and proofs of $10^2$\,MB per cube, i.e.\ $10^4$--$10^5$ CPU-hours for
all of $n=18$---feasible on a cluster but not done.  (ii) Proof logging inside
the forest engine (VeriPB/pseudo-Boolean), which would certify the isomorph
rejection as well; not attempted.  We regard the present level of verification
(independent oracle for the shallow levels, differential tests on planted
instances, three bit-identical replications, and an independent implementation
reproducing all per-level counts) as adequate for the claim, and the per-level
counts as the reproducible artefact that a sceptical reader should target.

\paragraph{What remains.}
The next admissible orders are $n=25$ ($N=300$) and $n=27$ ($N=351$).
Extrapolating the growth factor of $6.5$--$7$ per unit of $n$ gives of the
order of $10^{16}$--$10^{17}$ nodes at $n=25$, which is beyond the present
method by four to five orders of magnitude; new pruning ideas (for instance
exploiting the top-value structure and the parity classes inside the forest
search, or a meet-in-the-middle over the smallest and largest weights) or a
genuinely different argument would be needed.  Whether Leech's five trees are
the only Leech trees remains open; the present note only moves the smallest open
order from $18$ to $25$.  In the neighbouring problems of
Section~\ref{sec:variants} the frontier is much closer: the exact value of
$M(12)$ needs one exhaustive run that is queued, $M(13)$ is out of reach of the
present engine, modular Leech trees of order $12$ and above are untouched, and
the leaf-Leech question at $9$ leaves needs an engine of the kind built here for
Leech trees rather than a generic solver.

\section{AI-assisted research statement}\label{sec:ai}

The search programs, the test and verification harnesses, the literature
ledger and the drafts of this note were produced by AI coding agents (Claude,
Anthropic, in the Claude Code environment) working under the direction of the
author, who set the problem, chose the algorithmic route (forced-forest search
with isomorph rejection), decided what had to be verified, ran the cluster jobs
and reviewed the proofs and the code.  An adversarial ``referee'' pass over both
engines was likewise carried out by an AI agent instructed to look for unsound
pruning and incomplete symmetry breaking; it produced the correction to
Lemma~\ref{lem:sidon} recorded in Section~\ref{sec:prelim} and the verification
scheme of Section~\ref{sec:verification}.  The clean-room implementation of
Section~\ref{sec:cleanroom} was written by a separate agent that had no access to
the production source.  All mathematical claims in this note were checked by
independent implementations as described in Section~\ref{sec:verification};
human review of the proofs, the code and the text is ongoing, and the author
takes responsibility for the content.
\todo{author to confirm and adjust this statement}

\section{Data and code availability}\label{sec:data}

All code and data are available at
\url{https://github.com/lsmeng/leech-trees}
\todo{repository to be created; add archival DOI (Zenodo) at submission}.
The archive contains: the two search engines (\file{forest\_search.cpp},
\file{leech\_search.cpp}) and the clean-room \file{cr\_forest.c}; the
independent Python enumerators and differential harnesses; the SAT encoder and
the DRAT-verified proofs for $n\le 11$; the frozen list of the $123{,}867$ trees
of order $18$ with the two-generator cross-check; the per-shard output records
and depth histograms of the $n=18$ runs and of the $n\le 17$ runs; the verdict
scripts; and the referee reports on which Section~\ref{sec:verification} is
based.  It also contains, in a separate directory, the modified engines and the
raw outputs behind Section~\ref{sec:variants} (minimal distinct distance trees,
modular Leech trees, leaf-Leech trees), including the witnesses printed above and
the independent checker used on them.  The five known Leech trees are checked by
two independent checkers, and every witness produced by any engine in the project
was passed through both.

\subsection*{Acknowledgements}
Computations were carried out on the author's laptop and on the Hoffman2
Shared Cluster of the UCLA Office of Advanced Research Computing.
\todo{confirm acknowledgement wording; add funding if any}

\bibliographystyle{abbrvnat}
\bibliography{refs}

\end{document}
